\section{The full ideal range and its sharpness}
\label{sec:ideal}

\begin{proof}[Proof of Theorem~\ref{thm:main}]
The first two assertions are Corollary~\ref{cor:tails} and
Proposition~\ref{prop:operator-convergence}. It remains to prove
the ideal convergence and its exact range.

Fix a real $q>0$ with $q\rho>1$. Choose
\begin{equation}
 \frac1q<\eta<\rho,
 \label{eq:eta-choice}
\end{equation}
and let $M$ be a common constant for the two bounds in
Corollary~\ref{cor:tails}. Put $R_N=A_N-E_\sigma$. By
Lemma~\ref{lem:sv-sum} and positivity of $A_N,E_\sigma$,
\begin{equation}
 s_{2j-1}(R_N)
 \le s_j(A_N)+s_j(E_\sigma)
 \le2M j^{-\eta}
 \qquad(j\ge1).
 \label{eq:difference-tail}
\end{equation}
The even singular value $s_{2j}(R_N)$ is no larger than the preceding
one. Consequently,
\begin{equation}
 \sum_{k\ge1}s_k(R_N)^q
 \le2(2M)^q\sum_{j\ge1}j^{-\eta q}<\infty,
 \label{eq:summable-majorant}
\end{equation}
uniformly in $N$. In particular, every $R_N$ belongs to $\Sq$
in this range. For each fixed $k$,
$s_k(R_N)\le\opnorm{R_N}\to0$ by
Proposition~\ref{prop:operator-convergence}. The majorant in
\eqref{eq:difference-tail}, grouped in pairs of indices, is
$q$-summable because $\eta q>1$. Dominated convergence for
counting measure on $\N$ now gives
\[
 \qnorm{A_N-E_\sigma}^q
 =\sum_{k\ge1}s_k(R_N)^q\longrightarrow0.
\]
This proof uses scalar summation of singular values; no Banach-norm
interpolation or triangle inequality is assumed when $q<1$.

Finally suppose $q\rho\le1$. By
Theorem~\ref{thm:classical}, $E_\sigma\notin\Sq$. For every
finite $N$, however, $A_N$ is finite rank and hence belongs to
$\Sq$. If $A_N-E_\sigma$ belonged to $\Sq$, closure under
differences from Lemma~\ref{lem:sv-sum} would imply
$E_\sigma=A_N-(A_N-E_\sigma)\in\Sq$, a contradiction.
This proves the stated nonmembership for every $N$, and completes
the theorem.
\end{proof}

The negative conclusion is stronger than failure of a uniform bound
or failure of a chosen convergence argument. No finite Gram matrix can
remove the non-summable tail of the classical limiting operator.
Conversely, the positive proof requires no endpoint estimate for
the approximating eigenvalues: the choice~\eqref{eq:eta-choice}
uses only an exponent strictly between the ideal threshold and
$\rho$.
